\subsection{Nilai-Nilai Luhur dalam Perumusan dan Pengesahan Dasar Negara}
\begin{enumerate}
\item Menghargai pendapat.
\item Rendah hati.
\item Kerja keras.
\item Keberanian.
\item Rela berkorban.
\item Menerima keputusan rapat.
\item Mengutamakan : Persatuan \& kesatuan.
\item Melaksanakan keputusan rapat.
\end{enumerate}

\subsection{Dinamika Dasar Negara Pancasila}
UUD 1945 mengalami beberapa perubahan
\begin{table}[H]
  \centering
  \begin{tabular}{ccc}
    \toprule
    \textbf{Sila} & \textbf{Konstitusi RIS} & \textbf{UUD 1950}\\ \midrule
    Pertama & Ketuhanan YME & Ketuhanan YME \\
    Kedua & Perikemanusiaan & Kemanusiaan yang adil dan beradab \\
    Ketiga & Kebangsaan & Persatuan Indonesia \\
    Keempat & Kerakyatan & Kedaulatan rakyat \\
    Kelima & Keadilan sosial & Keadilan sosial \\ \bottomrule
  \end{tabular}
\end{table}

Sila-sila pancasila diperkuat dengan instruksi presiden RI No 12 tahun 1968 / 13 April 1968.


Setelah keluar instruksi presiden, adanya penyimpangan terhadap UUD $\rightarrow$ Ancaman dari ideologi komunis. Kudeta yang sebelumnya dinamakan operasi Takari $\rightarrow$ Gerakan 30 September $\rightarrow$ Mereka menculik dan membunuh para perwira tinggi A.D.


Yang gugur $\rightarrow$ Letnah Jendral Ahmad Yani, Mayor Jendral Raden Soeprapto, Mayor Jendral Mas Tirtodarmoharyono, Brigadir Jendral Siswando, Parman, dan Letnan Satu Dierre Andreas (Ajudan Kasab Jendral Nasution).


PKI dibubarkan $\rightarrow$ Ketetapan MPRS nomor XXV/MPRS/1966 tentang pembubaran partai komunis Indonesia.


Tanggal 1 Oktober ditetapkan $\rightarrow$ Hari kesaktian pancasila. Berdasarkan keputusan presiden No 153 tahun 1967 dan lubang buaya dijadikan monumen pancasila sakti.


\subsection{Alur perubahan UUD}
UUD (1945) $\rightarrow$ UU RIS (1949 - 1950) $\rightarrow$ UUDS 1950 (1950 - 1959) $\rightarrow$ Amandemen UUD 1945 (1959 - Sekarang).
